\section{A quartic equation in forty-two unrestricted integers}
\label{sec:quartic42}

The new equation in this section has \emph{forty-two unrestricted integer
unknowns}, exact total degree four, and twenty-seven quadratic residuals.
It improves the supplied forty-five-integer quartic without increasing its
degree. The four fixed coefficient ports are the complete, unchanged tuple
of that predecessor. In particular, the high-mask port is already shifted.
The construction uses the direct Pell/compiler route; no Jones or Sun
variable-reduction theorem is needed for this result.

\paragraph{Theorem.}
With the fixed coefficients specified in the accompanying binding, the
polynomial $Q_{42}$ defined below satisfies
\[
 \mathrm{ZFC}\ \vdash\
 \bigl[\neg\exists\boldsymbol{x}\in\mathbb Z^{42}
                \;Q_{42}(\boldsymbol{x})=0\bigr]
 \ \longleftrightarrow\ \operatorname{Con}(\mathrm{ZFC}).
\]
The argument is an ordinary arithmetic proof, composed with the expressly
retained native compiler and machine-calibration theorems. The new literal
polynomial has not received an end-to-end Lean or Rocq certificate. An
exported derivation in the repository's object-language ZFC calculus is
also distinct from the mathematical proof given here.

\subsection{The idea: prove the sign after recovering the rank}

Two positive coordinates in the predecessor, $\eta$ and $\zeta$, are each
represented by a ternary positive form. Those two representations use six
extra integer unknowns. We replace them by one representation of
$\eta\zeta>0$, which uses three. This initially proves only that the two
coordinates have the same sign.

An existing square scale supplies the missing information. Replacing
$w=u^2$ by $w=4u^2$ makes a modulus $e$ more than four times the decoded
index $R$. The Pell equations first recover a rank $p=R$ without deciding
the common sign. A negative common sign would then require a positive
integer smaller than $e$ to be divisible by $e$. Thus both coordinates
are positive after all. Canonical accepting witnesses have $w$ equal to
a sufficiently large dyadic square, so the factor four does not obstruct
completeness.

This ordering is essential. Invoking the predecessor's whole soundness
theorem before recovering the missing positive guards would be circular.
The local rank argument below is what permits those guards to be removed.

\subsection{Every variable and every residual}

Write the four fixed positive integer ports as
$\beta,K_0,M_C,M_F$, where $\beta=2^{1064921}-1$ and $M_F$ denotes the
\emph{already shifted} source mask. The complete coefficient binding, not
that exponent alone, specifies the instance. Let the thirty-three
arithmetic unknowns be
\begin{gather*}
 J,F,\alpha,z,f,h,r,v,w,\tau,\eta,\zeta,y,\gamma,o,j,\\
 N,E,P,c,a,X,D,C_2,L,F_2,V,g,I,B_*,\Delta,u,\rho.
\end{gather*}
Add nine guard unknowns $A_1,A_2,A_3,B_1,B_2,B_3,T_1,T_2,T_3$.
All forty-two range over $\mathbb Z$. In particular, $C_2,F_2,B_*$ and
$\Delta$ are independent supplied coordinates; their defining equations
are among the residuals. Define only the following abbreviations:
\[
 q=\beta J+1,\qquad C=q-F-\alpha,\qquad k=\eta+\zeta,
 \qquad \operatorname{Pos}(x,y,t)=1+x^2+y^2+t(t+1).
\]
These abbreviations add no quantified unknowns. The residuals are
\begin{align}
 R_1&=L-q^2+1, &
 R_2&=N-q(L+1), \nonumber\\
 R_3&=X-wq, &
 R_4&=a-E-(2v+1)N, \nonumber\\
 R_5&=E(X+1)-Xa, &
 R_6&=c-\eta-(a-E)k, \nonumber\\
 R_7&=P-k-E(c-\eta), &
 R_8&=\tau^2-P(P-k)-1, \nonumber\\
 R_9&=D-X-ac-\gamma(4a+3). &&
 \label{eq:q42-block1}
\end{align}
The next block is
\begin{align}
 R_{10}&=\Delta-a^2-4a-3, \nonumber\\
 R_{11}&=C_2-c^2, \nonumber\\
 R_{12}&=D^2-\Delta C_2-1, \nonumber\\
 R_{13}&=k-1-hE
       -\bigl[(q-F)(N-q)-CL+(M_C+qM_F)J\bigr], \nonumber\\
 R_{14}&=(K_0+w)C-F-z(q-1), \nonumber\\
 R_{15}&=F_2-f^2, \nonumber\\
 R_{16}&=V+c-fo, \nonumber\\
 R_{17}&=V-jc+k-1-hE, \nonumber\\
 R_{18}&=g+V^2-y^2.
 \label{eq:q42-block2}
\end{align}
Finally, put
\begin{align}
 R_{19}&=I-(2r+1)C_2, &
 R_{20}&=B_*-I^2, \nonumber\\
 R_{21}&=B_*-\Delta(F_2-1), &
 R_{22}&=y^2-B_*g-1, \nonumber\\
 R_{23}&=w-4u^2, &
 R_{24}&=q-2\rho^2, \nonumber\\
 R_{25}&=F\alpha-\operatorname{Pos}(A_1,A_2,A_3), &&\nonumber\\
 R_{26}&=FC-\operatorname{Pos}(B_1,B_2,B_3), &&\nonumber\\
 R_{27}&=\eta\zeta-\operatorname{Pos}(T_1,T_2,T_3). &&
 \label{eq:q42-block3}
\end{align}
The required single equation is
\begin{equation}
 \boxed{\displaystyle Q_{42}=\sum_{i=1}^{27}R_i^2=0.}
 \label{eq:q42-polynomial}
\end{equation}
Every residual has degree at most two in the supplied unknowns. The
coordinate $\tau$ occurs only in $R_8$, with term $\tau^2$, so the
coefficient of $\tau^4$ in $Q_{42}$ is exactly one. This proves exact
degree four on the authentic fixed-coefficient slice, independently of
any generic degree-propagation estimate.

The literal JSON uses \texttt{v\_scale}, \texttt{Baux} and
\texttt{q\_half\_root} for $v,B_*,\rho$, and named triples for the three
positive forms. The verifier expands all twenty-seven formulas above
and checks exact equality with that JSON. In particular, no division,
variable exponent or unspecified coefficient appears in the definition.

\subsection{Preliminary guards and restored arithmetic}

For every integer triple, $\operatorname{Pos}(x,y,t)>0$, since
$t(t+1)\ge0$. Conversely every positive integer $n$ has such a
representation. Indeed $4n-3\equiv1\pmod4$ is not of the forbidden form
$4^a(8b+7)$, so the three-square theorem represents it by three squares.
Exactly one square is odd. Writing that square as $(2t+1)^2$ and the
others as $(2x)^2,(2y)^2$ gives the claimed representation.

At an integer zero of \eqref{eq:q42-polynomial}, every residual vanishes.
The dyadic guard gives
\[
 q=\beta J+1=2\rho^2.
\]
Since $\beta\ge3$, $J<0$ would make $q<0$, while $J=0$ would require
$1=2\rho^2$. Thus $J>0$ and, for the bound source tuple, $q\ge16$.
The other two retained guards give $F\alpha>0$ and $FC>0$. Hence
$F,\alpha,C$ have the same nonzero sign; their sum is $q>0$, so
\begin{equation}
 F>0,\qquad \alpha>0,\qquad C>0.
 \label{eq:q42-simplex}
\end{equation}

Put $s=2v+1$, a nonzero odd integer, and set $Y=sq^3$ as a
\emph{proof abbreviation}. The first residuals give
\begin{gather*}
 N=q^3,\quad X=wq,\quad a=E+Y,\quad E=XY,\quad
 c=kY+\eta,\quad P=XY^2k+k,\\
 H=4a+3,\quad \Delta=(a+2)^2-1,\quad
 D=X+ac+\gamma H,\quad I=(2r+1)c^2.
\end{gather*}
For example $R_5=0$ and $a=E+Y$ imply
$E(X+1)=X(E+Y)$, hence $E=XY$ with no division or nonzero assumption.
Define the proof-only packing integer
\begin{equation}
 R=\bigl[(q-1)(q-F)+\alpha\bigr](q^2-1)
       +(M_C+qM_F)J.
 \label{eq:q42-packing}
\end{equation}
The retained equations then assert
\begin{gather}
 \tau^2-XY^2(XY^2+1)k^2=1,
 \qquad D^2-\Delta c^2=1, \label{eq:q42-pell1}\\
 I^2=\Delta(f^2-1),
 \qquad I^2V^2-(I^2-1)y^2=1, \label{eq:q42-pell2}\\
 k=R+1+hE,
 \qquad V=fo-c=jc-R. \label{eq:q42-index}
\end{gather}
The last two identities expose the quotient congruences directly; an
unproved quotient-integrality assertion is unnecessary.

The source packing bounds require only \eqref{eq:q42-simplex} and the
authentic mask bounds at this stage. With
$G=q\alpha+(q-1)C$, one has
\[
 q\le G\le q^2-q,
 \qquad 0<(M_C+qM_F)J<(q-1)(1+2q).
\]
Consequently
\begin{equation}
 q(q^2-1)<R<q^4-2,
 \qquad R-2>3q+1.
 \label{eq:q42-rbounds}
\end{equation}
No positivity assumption on $z,\eta,\zeta$ or $Y$ has been used in
obtaining these bounds.

\subsection{Bootstrap with an unknown common sign}

The new guard says $\eta\zeta>0$. Let
\[
 \sigma=\operatorname{sign}(\eta)=\operatorname{sign}(\zeta),
 \qquad \kappa=|k|>0,
 \qquad \theta=\frac{\eta}{k}\in(0,1).
\]
The fraction $\theta$ is used only in the proof. First $u\ne0$.
Otherwise $w=X=E=0$, $a=Y$ and $D=ac+\gamma H$. Reducing the main
Pell norm modulo $H$ would give $0\equiv1\pmod H$, since
$\Delta\equiv a^2\pmod H$. But $|Y|\ge q^3$ makes
$|H|=|4Y+3|>1$. This is impossible. Therefore
\begin{equation}
 w=4u^2\ge4,\quad X\ge4q>0,\quad |Y|\ge q^3,
 \quad e:=|E|\ge4q^4>4R+8.
 \label{eq:q42-modulus}
\end{equation}

Set $A=|a+2|$, $P_0=2XY^2+1$, and $c_0=|c|$. Then
\[
 A>e>R+2,\qquad P_0>A\ge2,\qquad \Delta=A^2-1.
\]
For negative $Y$, the first inequality follows from
$A=(X+1)|Y|-2>e$; for positive $Y$ it is immediate. Moreover
\[
 \frac{c_0}{\kappa}=
 \begin{cases}
 |Y|+\theta,&Y>0,\\
 |Y|-\theta,&Y<0,
 \end{cases}
\]
so $c_0>\kappa>0$ in either case.

Let $\chi_A(n),\psi_A(n)$ be the integer Pell sequences defined by
\[
 \chi_A(n)+\psi_A(n)\sqrt{A^2-1}
       =(A+\sqrt{A^2-1})^n.
\]
Ordinary Pell classification in \eqref{eq:q42-pell1} gives positive
ranks $n,p$ with
\begin{equation}
 \kappa=2\psi_{P_0}(n),\quad |\tau|=\chi_{P_0}(n),
 \qquad c_0=\psi_A(p),\quad |D|=\chi_A(p).
 \label{eq:q42-ranks}
\end{equation}
The first fundamental unit is the standard
$2b+1+2\sqrt{b(b+1)}$, where $b=XY^2>1$.
Since $P_0>A$ and $c_0>\kappa$, strict Pell growth gives $p>n$.
Also $P_0\equiv1\pmod e$, so the recurrence gives
\begin{equation}
 \sigma\,2n\equiv R+1\pmod e.
 \label{eq:q42-signedindex}
\end{equation}
If $\sigma=1$, its least positive representative gives $2n\ge R+1$.
If $\sigma=-1$, it gives $2n\ge e-R-1>R+1$. Thus
\[
 n\ge(R+1)/2\ge24,\qquad p>n.
\]
In particular the elementary Pell bounds imply
\begin{equation}
 c_0>A\Delta^2,\qquad c_0>2p,\qquad c_0>2R.
 \label{eq:q42-cbounds}
\end{equation}
For the first, $p\ge6$ already suffices, since
$\psi_A(6)=32A^5-32A^3+6A>A(A^2-1)^2$. The second follows by the
recurrence for $A\ge2,p\ge3$; the third follows from
$c_0\ge\psi_A(2)=2A>2R$.

\subsection{Recovering the rank without a sign assumption on \texorpdfstring{$c$}{c}}

We use the coefficient-independent ordinary strong-rank lemma retained
with the source proof. In its needed form, if
\[
 c_0=\psi_A(p)>A(A^2-1)^2,
 \qquad (i c_0^2)^2=(A^2-1)(f^2-1),\quad i>0,
\]
then there is a positive integer $m$ such that
\begin{equation}
 |f|=\chi_A(m),\quad p\mid m,\quad c_0\mid m,
 \quad m\ge c_0>2p,
 \quad i c_0^2=(A^2-1)\psi_A(m).
 \label{eq:q42-strong}
\end{equation}
Here its hypotheses hold with $i=|2r+1|$, using
\eqref{eq:q42-pell2} and \eqref{eq:q42-cbounds}.

For clarity, the two substantive ingredients in that inherited lemma are
strong divisibility and an elementary growth comparison. Write the
quadratic field's fundamental integer-coordinate unit using a parameter
$F_0$, and write $A=\chi_{F_0}(d)$,
$S=\psi_{F_0}(d)$, and
$\Delta=(F_0^2-1)S^2$. Classification yields
$|f|=\chi_{F_0}(t)$ and
$|I|=(F_0^2-1)S\psi_{F_0}(t)$.
Strong divisibility, together with $c_0^2\mid I$, gives
$\psi_{F_0}(\gcd(t,dp))\ge c_0/\Delta$.
If that gcd were a proper divisor of $dp$, doubling its index would give
\[
 2\psi_{F_0}(\gcd(t,dp))^2
   <\psi_{F_0}(2\gcd(t,dp))
   \le S c_0<A c_0,
\]
contrary to $c_0>A\Delta^2$. Hence $dp\mid t$, giving
$t=dm$ and $p\mid m$. Writing $m=pz$, the recurrence modulo $c_0^2$
gives $c_0\mid\Delta z$. Finally
$\gcd(c_0,\Delta)=\gcd(p,\Delta)$ follows from
$\psi_A(p)\equiv pA^{p-1}\pmod\Delta$, and gives $c_0\mid pz=m$.
The archived ordinary-kernel proof supplies the complete signed and
fundamental-unit details of this lemma.

The auxiliary norm in \eqref{eq:q42-pell2} excludes $V=0$ and $y=0$.
Pell classification at the integer parameter $|I|>1$ gives an odd rank
$\ell=2v_0+1$ such that
\[
 |I|\,|V|=\chi_{|I|}(\ell),\qquad
 |y|=\psi_{|I|}(\ell),\qquad |V|=Q_{v_0}(I^2).
\]
Oddness follows from the even-index residue
$\chi_{|I|}(2t)\equiv(-1)^t\pmod{|I|}$.
The standard odd-subsequence polynomials satisfy
\[
 Q_{v_0}(1-A^2)=(-1)^{v_0}\psi_A(\ell),
 \qquad Q_{v_0}(0)=(-1)^{v_0}\ell.
\]
They follow directly from $Q_0(x)=1$, $Q_1(x)=4x-3$ and
$Q_{t+1}(x)=(4x-2)Q_t(x)-Q_{t-1}(x)$.

The exposed identities in \eqref{eq:q42-index} imply
\[
 V\equiv-c\pmod{|f|},\qquad V\equiv-R\pmod{c_0}.
\]
As $I^2\equiv1-A^2\pmod{|f|}$, the first congruence gives
$\psi_A(\ell)\equiv\pm c_0\pmod{|f|}$, with all signs of $c,V$
absorbed in the ambiguity. Squaring and applying the Pell doubling formula,
\[
 \chi_A(2\ell)\equiv\chi_A(2p)\pmod{\chi_A(m)}.
\]
The signed step-down argument gives $\ell\equiv\pm p\pmod m$.
Indeed reduce $2\ell$ modulo $2m$, using
$\chi_A(t+2m)\equiv-\chi_A(t)$, then reflect indices above $m$.
The resulting representative is $\pm\chi_A(s)$, $0\le s\le m$.
The case $s=m$ gives residue zero. A negative representative is impossible
because
$\chi_A(s)+\chi_A(2p)\le2\chi_A(m-1)<\chi_A(m)$.
Strict growth then forces $s=2p$, proving the claimed index congruence.
Its hypotheses $0<2p<m$ were proved in \eqref{eq:q42-strong}.

Since $c_0\mid m$, this gives $\ell\equiv\pm p\pmod{c_0}$.
On the other hand $c_0\mid I$, so the second exposed congruence and
$Q_{v_0}(0)=(-1)^{v_0}\ell$ give
$\ell\equiv\pm R\pmod{c_0}$. Therefore
$p\equiv\pm R\pmod{c_0}$. The strict bounds
$0<p,R<c_0/2$ in \eqref{eq:q42-cbounds} force
\begin{equation}
 p=R.
 \label{eq:q42-mainrank}
\end{equation}

\subsection{The missing signs and the reverse construction}

If the common sign were negative, \eqref{eq:q42-signedindex} would make
$e$ divide $2n+R+1$. But $n<p=R$ and \eqref{eq:q42-modulus} give
\[
 0<2n+R+1<3R+1<e,
\]
a contradiction. Thus $\eta,\zeta>0$. Represent them separately by
$\operatorname{Pos}$, and set the old square coordinate to
$u_{\mathrm{old}}=2u$. Every residual of the authenticated predecessor
now vanishes. Its remaining signed-scale and compiler arguments therefore
produce a finite ZFC refutation. The missing individual guards have been
restored before invoking that whole-system theorem.

There is a small transcription correction in the predecessor's written
transport discussion. Its literal circuit, retained in $R_{14}$, has
\[
 (K_0+w)C=F+z(q-1),
\]
whereas the prose writes $z-1$. The needed conclusion is $z>0$.
After its Pell recovery, $w>q>F$ and $C\ge1$ make the left side greater
than $F$, hence $z(q-1)>0$. This proves the required guard for the
literal shifted coordinate. Neither the equation nor its coefficients
are changed by correcting that sentence.

Conversely, a finite refutation supplies the predecessor's canonical
accepting witnesses. Its complete coefficient adapter ensures
\[
 q=2^t,\qquad t\ \hbox{odd},\qquad R\ \hbox{odd},
 \qquad w=2^{R-t},\qquad R>3q.
\]
Thus $R-t$ is even and at least four. Choose
\[
 u=2^{(R-t-2)/2},
\]
so $w=4u^2$. The canonical $\eta,\zeta$ are positive; represent their
product by one $\operatorname{Pos}$ triple and keep the other coordinates.
All twenty-seven new residuals vanish. This is an existence construction
using canonical witnesses, rather than an invertible substitution on
every predecessor tuple.

\subsection{What has been checked}

The new construction retains twenty-five predecessor residuals exactly,
strengthens the square scale, and replaces two positive guards by one.
It changes $(45,28,4)$ to $(42,27,4)$ in the order
unknowns, residuals, and single-polynomial degree.
With the four ports left symbolic, the exact expanded sum of squares
has 391 terms, compared with 407 for the predecessor. This is a symbolic
description count, not a claim about coefficient bit complexity or an
optimal arithmetic circuit.

The main builder and a separate implementation expand the literal
residuals in sparse integer polynomial arithmetic. They check variable
counts and occurrence, the unchanged rows, the two modifications, exact
degree four, and the coefficient one of $\tau^4$. Both work under
ordinary and optimized Python. A separate transcription check verifies
all twenty-seven factored formulas above against the emitted JSON.
The parent polynomial and its proof are authenticated by their pinned
Git blob hashes; the complete coefficient tuple and semantic source
dependencies are retained with their own provenance.

These finite checks establish the displayed algebra and source binding.
The universal solvability equivalence rests on the two-way mathematical
argument and its expressly identified inherited arithmetic and compiler
theorems. No consistency assumption, refutation search, global
minimality assertion or external peer review is being substituted for
that proof.
